\documentclass[12pt,a4paper]{article}
% Estilo compartido para informes. Compilar con LuaLaTeX desde tp1/.
\usepackage[spanish,es-nodecimaldot]{babel}
\usepackage{fontspec}
\setmainfont{Montserrat-Medium.ttf}[
  Path=../assets/fonts/,
  BoldFont=Montserrat-Bold.ttf,
  ItalicFont=Montserrat-MediumItalic.ttf]
\usepackage[a4paper,margin=2.5cm,headheight=16pt]{geometry}
\usepackage{xcolor,graphicx,booktabs,tabularx}
\usepackage{float}
\usepackage{titlesec,fancyhdr,enumitem,caption}
\usepackage{ragged2e}
\usepackage{hyperref}
\definecolor{verde}{HTML}{1D4036}
\definecolor{bronce}{HTML}{D38129}
\definecolor{rosa}{HTML}{D88E77}
\definecolor{calido}{HTML}{EFE4D8}
\definecolor{grafito}{HTML}{1D1D1D}
\hypersetup{colorlinks=true,linkcolor=verde,urlcolor=verde,
  pdftitle={TP1 - Instalación de máquinas virtuales},pdfsubject={Cyberseguridad}}
\AtBeginDocument{\color{grafito}\RaggedRight\fontsize{12}{14.4}\selectfont}
\setlength{\parindent}{0pt}
\setlength{\parskip}{7pt}
\setlist[itemize]{leftmargin=*,itemsep=7pt,topsep=4pt}
\setlist[enumerate]{leftmargin=*,itemsep=6pt}
\titleformat{\section}{\color{verde}\bfseries\fontsize{18}{21.6}\selectfont}{\thesection}{0.65em}{}
\titleformat{\subsection}{\color{verde}\bfseries\large}{\thesubsection}{0.65em}{}
\titleformat{\subsubsection}{\color{verde}\bfseries}{\thesubsubsection}{0.65em}{}
\titlespacing*{\section}{0pt}{16pt}{8pt}
\titlespacing*{\subsection}{0pt}{12pt}{5pt}
\titlespacing*{\subsubsection}{0pt}{10pt}{4pt}
\setcounter{tocdepth}{2}
\pagestyle{fancy}
\fancyhf{}
\fancyhead[L]{\small\color{verde}FCEFyN | UNC}
\fancyhead[R]{\small\color{verde}\materia\ · TP \numeroTP}
\fancyfoot[L]{\scriptsize Informe de laboratorio}
\fancyfoot[R]{\small\thepage}
\renewcommand{\headrulewidth}{0.4pt}
\captionsetup{font=small,labelfont=bf,justification=raggedright,singlelinecheck=false,hypcap=false}
\newcommand{\pendiente}[1]{\par{\small\color{verde}\textbf{Por completar.} #1\par}}
% Si falta una captura se muestra un espacio reservado; no impide compilar.
\newcommand{\captura}[2]{%
  \par\smallskip\begingroup\centering
  \IfFileExists{figuras/#1}{%
    \includegraphics[width=\linewidth,height=5cm,keepaspectratio]{figuras/#1}%
  }{%
    \setlength{\fboxsep}{10pt}%
    \fcolorbox{bronce}{calido}{\parbox[c][1.25cm][c]{\dimexpr\linewidth-2\fboxsep-2\fboxrule\relax}{%
      \centering\small Captura pendiente\par\texttt{\detokenize{#1}}}}%
  }%
  \captionof{figure}{#2}\par\endgroup\smallskip
}

\usepackage{xurl}

\newcommand{\materia}{Cyberseguridad}
\newcommand{\numeroTP}{17}
\newcommand{\tituloTP}{Investigación de incidentes de seguridad}
\newcommand{\estudiante}{Brigido Tomas Andres}
\newcommand{\legajo}{38481096}
\newcommand{\carrera}{ICOMP}
\newcommand{\docente}{Javier A. Jorge}


% Metadatos específicos del TP17 y capturas en img/.
\hypersetup{pdftitle={TP17 - Investigación de incidentes de seguridad},linkcolor=black,urlcolor=black,citecolor=black}
\renewcommand{\captura}[2]{%
  \begin{figure}[H]
    \centering
    \IfFileExists{img/#1}{%
      \includegraphics[width=0.85\linewidth,height=6cm,keepaspectratio]{img/#1}%
    }{%
      \setlength{\fboxsep}{10pt}%
      \fcolorbox{bronce}{calido}{\parbox[c][1.1cm][c]{0.85\linewidth}{%
        \centering\small Captura pendiente\par\texttt{\detokenize{#1}}}}%
    }%
    \caption{#2}
  \end{figure}
}

\usepackage{caption}
\captionsetup{justification=centering}
\begin{document}
\hypersetup{pageanchor=false}
\begin{titlepage}
  \includegraphics[width=8cm]{../assets/logo-fcefyn-2026.png}
  \vspace{1cm}

  {\color{verde}\small UNIVERSIDAD NACIONAL DE CÓRDOBA\par}
  \vspace{2.1cm}
  {\color{bronce}\bfseries\large TRABAJO PRÁCTICO \numeroTP\par}
  \vspace{0.45cm}
  {\color{verde}\bfseries\fontsize{32}{38.4}\selectfont Investigación de\\incidentes de seguridad\par}
  \vspace{0.3cm}
{\small Análisis de registros con Security Onion\par}
\vspace{0.35cm}
  {\large \materia\par}
  \vspace{0.8cm}
  {\color{bronce}\rule{2.7cm}{2pt}\par}
  \vfill
  \renewcommand{\arraystretch}{1.45}
  \begin{tabular}{@{}p{3.2cm}p{10cm}@{}}
    Estudiante & \estudiante \\
    Legajo & \legajo \\
    Carrera & \carrera \\
    Docente & \docente \\
  \end{tabular}
  \vfill
  {\small Córdoba, Argentina\par}
\end{titlepage}
\hypersetup{pageanchor=true}

\tableofcontents
\vspace{1cm}
\newpage


% Cada pregunta deja un espacio editable para una respuesta breve.
\newcommand{\pregunta}[1]{\par{\sloppy\textbf{#1}\par}\nopagebreak\pendiente{Respuesta.}}
\section{Introducción}
Este trabajo reúne tres laboratorios de análisis de incidentes con Security Onion. Se utilizarán Kibana, Sguil y Wireshark para investigar tráfico HTTP y DNS, identificar comunicaciones mediante cinco tuplas y reconstruir un ataque a un equipo Windows.

\section{Laboratorio 27.2.12: análisis HTTP y DNS}
\subsection{Preparación y selección del período}
Comprobar los servicios con \texttt{sudo so-status} y seleccionar junio de 2020 en Kibana.
\captura{so-status.png}{Estado de los servicios de Security Onion. Logstash muestra una advertencia de inicialización.}
\captura{kibanaJune.png}{Panel general de Kibana y registros del período seleccionado.}

\clearpage
\subsection{Investigación de la inyección SQL}
Filtrar el tráfico HTTP y revisar los detalles de las solicitudes.
\captura{httplogs.png}{Registros HTTP en Kibana.}
\par{\sloppy\textbf{¿Cuál es la dirección IP de origen?}\par}
\texttt{209.165.200.227}.

\par{\sloppy\textbf{¿Cuál es la dirección IP de destino?}\par}
\texttt{209.165.200.235}.

\par{\sloppy\textbf{¿Cuál es el número de puerto de destino?}\par}
\texttt{80} (HTTP).

\clearpage
\captura{uriField.png}{Detalle del mensaje HTTP y del campo URI.}
\par{\sloppy\textbf{¿Cuál es la marca de tiempo del primer resultado?}\par}
12 de junio de 2020, a las 21:30:09.445, según la hora mostrada en Kibana.

\par{\sloppy\textbf{¿Cuál es el tipo de evento?}\par}
\texttt{bro\_http}, correspondiente a un registro HTTP.

\par{\sloppy\textbf{¿Qué se incluye en el campo de mensaje? Enfocarse especialmente en el campo URI.}\par}
Incluye los detalles de una solicitud GET. El URI contiene una inyección \texttt{UNION SELECT} en el parámetro \texttt{username}, que consulta números de tarjeta, códigos de verificación y vencimientos de la tabla \texttt{credit\_cards}.

\par{\sloppy\textbf{¿Qué es importante sobre esta información?}\par}
Evidencia un intento de extraer datos sensibles mediante inyección SQL. El código HTTP 200 por sí solo no confirma la extracción.


\clearpage
\subsection{Revisión de la transcripción en capME!}
\pendiente{Agregar una captura de la transcripción y de los datos identificados.}
\par{\sloppy\textbf{¿Qué puede ver más adelante en la transcripción en cuanto a nombres de usuario?}\par}
En la transcripción de la consigna aparecen números de tarjetas bajo la etiqueta Username, códigos de verificación como Password y fechas de vencimiento como Signature. No son credenciales de usuarios, sino datos de tarjetas extraídos.

\par{\sloppy\textbf{Dé algunos ejemplos de un nombre de usuario, contraseña y firma que se exfiltró.}\par}
Ejemplos de la captura de la consigna (página 5), usando sus etiquetas:
\begin{itemize}[itemsep=2pt]
    \item Username: \texttt{4444111122223333}; Password: \texttt{745}; Signature: \texttt{2012-03-01}.
    \item Username: \texttt{774653633776330}; Password: \texttt{722}; Signature: \texttt{2015-04-01}.
\end{itemize}


\clearpage
\subsection{Exfiltración mediante DNS}
Filtrar las consultas relacionadas con \texttt{example.com}, exportar los subdominios y decodificar las cadenas hexadecimales.
\pendiente{Agregar capturas del filtro DNS, las consultas y el texto decodificado.}
\pregunta{Registre las direcciones IP del cliente y servidor DNS.}
\pregunta{¿Los subdominios de las consultas DNS eran nombres legítimos? Si no, ¿cuál es el texto?}
\pregunta{¿Qué implica este resultado sobre esta respuesta DNS en particular? ¿Cuál es el mayor significado?}
\pregunta{¿Qué puede haber creado estas consultas DNS codificadas y por qué se seleccionó DNS como medio para exfiltrar datos?}

\clearpage
\section{Laboratorio 27.2.14: análisis de cinco tuplas}
Las cinco tuplas corresponden a las direcciones IP y los puertos de origen y destino, junto con el protocolo utilizado. Se investigará qué ocurrió con el archivo \texttt{confidential.txt}.
\subsection{Revisión de alertas en Sguil}
\pendiente{Agregar la alerta y su transcripción. Los ID pueden variar entre instalaciones.}
\pregunta{¿Qué tipo de transacciones ocurrieron entre el cliente y el servidor en este ataque?}

\subsection{Análisis del flujo TCP en Wireshark}
\pendiente{Agregar una captura de Follow TCP Stream.}
\pregunta{¿Qué observó? ¿Qué indican los colores de texto rojo y azul?}
\pregunta{El atacante emite el comando \texttt{whoami} en el objetivo. ¿Qué muestra esto sobre el rol de atacante en el equipo de destino?}
\pregunta{¿Qué tipo de datos ha estado leyendo el agente de amenaza?}

\clearpage
\subsection{Investigación de FTP y archivos en Kibana}
Seleccionar junio de 2020, incluyendo el día 11, y revisar los registros FTP y FTP\_DATA.
\pendiente{Agregar capturas de los registros, las cinco tuplas y el contenido transferido.}
\pregunta{¿Cuáles son las direcciones IP y los números de puerto de origen y de destino para el tráfico FTP?}
\pregunta{¿Cuáles son las credenciales de usuario para acceder al sitio FTP?}
\pregunta{¿Cuáles son los diferentes tipos de archivos? Mire la sección Tipo MIME.}
\pregunta{¿Cuáles son los orígenes de archivos enumerados?}
\pregunta{¿Cuál es el tipo MIME, la dirección IP de origen y destino asociado con la transferencia de los datos FTP? ¿Cuándo ocurrió esta transferencia?}
\pregunta{¿Cuál es el contenido de texto del archivo que se transfirió mediante FTP?}
\pregunta{Con toda la información recopilada hasta ahora, ¿cuál es su recomendación para detener nuevos accesos no autorizados?}

\clearpage
\section{Laboratorio 27.2.16: ataque a Windows}
Investigar las alertas del 19 de marzo de 2019 para identificar el equipo afectado, su usuario y el malware involucrado.
\subsection{Cronología y primera alerta en Sguil}
\pendiente{Agregar capturas de las alertas, los datos del paquete y la regla IDS.}
\pregunta{Según Sguil, ¿cuáles son las marcas de tiempo para la primera y última de las alertas que ocurrieron el 19/03/2019? ¿Qué es interesante acerca de las marcas de tiempo de todas las alertas de ese día?}
\pregunta{¿Cuál era la dirección IP de origen y el número de puerto y la dirección IP de destino y el número de puerto?}
\pregunta{¿Qué tipo de protocolo y solicitud o respuesta estuvo involucrado?}
\pregunta{¿Qué es la alerta y el mensaje de IDS?}
\pregunta{¿Cree que esta alerta fue el resultado de una mala configuración de IDS o de una comunicación sospechosa legítima?}
\pregunta{¿Cuál es el nombre de host, el nombre de dominio y la dirección IP del host de origen en la actualización DNS?}

\clearpage
\subsection{Primera descarga y análisis del archivo}
Exportar el objeto HTTP de la segunda alerta. Calcular su hash SHA256 y consultar su reputación.
\pendiente{Agregar la transcripción, el hash calculado y el resultado de la consulta.}
\pregunta{¿Cuáles son las direcciones MAC e IP y los números de puerto de origen y de destino?}
\pregunta{Mirando la solicitud (azul), ¿para qué fue la solicitud?}
\pregunta{¿Cuáles son los primeros caracteres del archivo? ¿Qué tipo de archivo identifica esa firma?}
\pregunta{¿Reconoció Talos el hash del archivo y lo identificó como malware? Si es así, ¿qué tipo de malware?}

\subsection{Comunicación asociada a Remcos}
\pendiente{Agregar la alerta y la transcripción correspondientes a Remcos.}
\pregunta{¿Cuál es el puerto de destino de la comunicación? ¿Es este un puerto conocido?}
\pregunta{¿Es legible la comunicación o está encriptada?}
\pregunta{¿Qué hace Remcos?}
\pregunta{¿Qué tipo de comunicación cree que se estaba transmitiendo?}
\pregunta{¿Qué tipo de encriptación y ofuscación se usó para eludir la detección?}

\clearpage
\subsection{Segunda descarga y correlación de alertas}
\pendiente{Agregar evidencia del segundo ejecutable y de su reputación.}
\pregunta{¿Qué ID alerta a un segundo archivo ejecutable que se está descargando?}
\pregunta{¿Desde qué dirección IP del servidor y número de puerto se descargó el archivo?}
\pregunta{¿Cuál es el nombre del archivo que se ha descargado?}
\pregunta{¿Es el archivo ejecutable malware conocido y, en caso afirmativo, de qué tipo? ¿Cuál es el nombre de la detección de AMP?}
\pregunta{¿Cómo se relacionan las tres alertas restantes?}

\subsection{Investigación complementaria en Kibana}
Seleccionar marzo de 2019 y acotar el intervalo al evento del día 19. Revisar alertas, HTTP y DNS.
\pendiente{Agregar capturas de las alertas y los paneles HTTP y DNS.}
\pregunta{¿Cuál es el país y la ciudad de la alerta asociada a 31.22.4.176?}
\pregunta{¿Cuál es el país y la ciudad de la alerta asociada a 115.112.43.81?}
\pregunta{¿Qué es el recuento de registros en el panel HTTP? ¿De qué países?}
\pregunta{¿Cuáles son los URI de los archivos descargados?}
\pregunta{¿A qué están relacionados los archivos \texttt{CSPCA.crl} y \texttt{ncsi.txt}?}
\pregunta{¿Alguno de los dominios parece potencialmente inseguro? ¿Cuál es el resultado de consultar \texttt{toptoptop1.online} en VirusTotal?}

\subsection{Síntesis del incidente}
\pregunta{¿Cuál fue la hora específica del ataque del 19/03/2019?}
\pregunta{¿Qué equipo host de Windows se infectó? ¿Quién era el usuario?}
\pregunta{¿Con qué estaba infectada la computadora?}

\clearpage
\section{Conclusión}
\pendiente{Resumir brevemente los hallazgos y la utilidad de correlacionar alertas, registros y tráfico.}
\end{document}
